\chapter*{怎样阅读本册}
\addcontentsline{toc}{chapter}{怎样阅读本册}

本册必须按顺序阅读，但顺序不是“名词清单”的顺序，而是一条失败推动机制出现的路线。先看一次普通保存请求为什么不能靠单个函数解释，再看硬件怎样提供执行现场、地址翻译、特权入口、中断和 DMA，随后才进入进程、调度、系统调用、虚拟内存、文件、设备、网络和安全。全书真实架构统一采用 x86-64，汇编统一采用 Intel 语法：目的操作数在前，源操作数在后，例如 \asmcode{mov rax, qword ptr [rbx]}、\asmcode{add rax, 1}、\code{syscall}、\code{iretq}。

第一部分回答“操作系统为什么必须存在”。先看 CPU、MMU、中断、DMA 和设备只暴露低层状态，再看裸机主循环为什么会被忙等、异步事件、死循环和启动约束击穿。读完这一部分，读者应该能把 task、address space、fd、request 这些 OS 对象追溯到具体机器事实。

第二部分回答“多个执行流怎样共享 CPU 且不丢事件”。它从进程、线程和 runqueue 进入，再补调度器、多核、等待队列、同步、锁和死锁。这里不要只背算法名，要看每个状态转换：谁从 runnable 变成 sleeping，谁负责唤醒，唤醒后怎样重新进入 runqueue。

第三部分回答“用户程序怎样被限制在授权边界内”。系统调用、\code{exec}、ELF、虚拟内存、缺页、\code{mmap}、page cache 和内核分配器被放在一起，是因为它们都在改变“谁能访问什么”。这一部分的主线是验证、准备、提交、回滚，避免半提交状态。

第四部分回答“异步设备怎样变成文件、网络和持久化语义”。驱动、DMA、块层、VFS、文件系统、socket、日志和恢复不是并列名词，而是一条从 request 身份、buffer 所有权、completion、page cache 到 \code{fsync} 的路径。

第五部分回答“系统出错、被攻击或性能异常时怎样给出证据”。安全隔离、panic、错误恢复、可观测性和验收方法统一收束到工程报告能力：结论必须能指向状态、边界、失败路径和可观察证据。

\begin{longtable}{@{}L{0.2\linewidth}L{0.34\linewidth}L{0.36\linewidth}@{}}
\toprule
阅读阶段 & 先问的问题 & 不合格读法 \\
\midrule
机器事实 & 哪个硬件限制逼出内核对象 & 只背 CPU、MMU、DMA 名词 \\
任务并发 & 状态怎样进入和离开 runqueue/wait queue & 只背调度算法复杂度 \\
边界内存 & 何时验证，何时提交，失败怎样回滚 & 把系统调用当普通函数 \\
I/O 持久化 & 请求身份、buffer 所有权和 completion 在哪 & 把 \code{write} 成功当落盘 \\
安全观测 & 哪条证据证明契约成立或失败 & 把工具输出直接当结论 \\
\bottomrule
\end{longtable}

学习每章时都要落到五个问题：

\begin{enumerate}
\tightlist
\item 状态在哪里：进程表、runqueue、锁等待队列、页表、VMA、fd 表、socket、inode、page cache、日志还是设备队列。
\item 硬件在哪里：RIP、RSP、RFLAGS、通用寄存器、cache、TLB、页表、MMIO 寄存器、DMA 描述符、APIC/MSI-X 还是定时器。
\item 权限在哪里：CPU 特权级、页表权限、文件权限、capability、namespace 还是对象引用。
\item 等待在哪里：runnable 等 CPU、blocked 等事件、mutex 等持有者、page fault 等页面、socket 等缓冲、I/O 等 completion。
\item 提交在哪里：\code{exec} 切换地址空间、状态更新完成、写入 page cache、块设备 flush、目录持久化还是 manifest 生效。
\end{enumerate}

每章可以按固定节奏读：先找失败场景，再看直接补丁为什么不够，再看 OS 机制如何建立状态对象，最后用不变量和证据收尾。若一节讲 fd，就要画出 fd table、open file description、inode 和 page cache 的关系；若一节讲调度，就要写出 runnable、running、blocked、zombie 和 reaped 的状态转换；若一节讲锁，就要画出锁顺序图和等待队列；若一节讲 \code{exec}，就要标出失败回滚边界和提交点；若一节讲持久化，就要写出每个崩溃点后恢复程序应该相信什么。

阅读实现小节时，不能只看伪代码能不能运行。要追问四件事：数据结构是什么，复杂度是多少，维护了哪些不变量，失败路径是否会留下半提交状态。一个教学模型可以比真实内核小，但不能把真实问题讲错。

\begin{warningbox}
不要把工具输出当成结论。工具只是证据来源；结论必须说明哪条 OS 契约被验证，哪条仍然只是合理假设。
\end{warningbox}
